%==============================================================================
% Section 4: P-channel dictionary and topology response
%==============================================================================
\section{P-channel dictionary and topology response}
\label{sec:pchannel}

\subsection{Channel separation}
\label{sec:pchannel:separation}

The Pontryagin-type quantities are separated into three channels:

\begin{table}[H]
\centering
\begingroup
\renewcommand{\arraystretch}{2.0}
\begin{tabular}{lL{4.5cm}L{3.5cm}L{3.5cm}}
\toprule
Channel & Result & Role & Distinction \\
\midrule
$\Pform$ & Exactly zero on active torsion branches
         & Structural cancellation
         & Distinct from dynamical / orbit stability \\
$\Pint$  & MX internal-pair diagnostic; $\Ctop=-9\chi$
         & Diagnostic topology response
         & Not a true Pontryagin density \\
$\QNY$   & Endpoint / boundary channel: zero on $\VT$,
           on-shell trivial on $\AX$, boundary-conditional on $\MX$
         & Boundary bookkeeping
         & Not a bulk dynamical mechanism \\
\bottomrule
\end{tabular}
\endgroup
\caption{P-channel dictionary.}
\label{tab:pchannel_dict}
\end{table}

$\Pint$ is a diagnostic quantity based on the contraction of internal index
pairs in the orthonormal frame, and is \emph{not} a form-Hodge Pontryagin
density~\cite{ChandiaZanelli1997}.  This distinction is essential for the
interpretation given in this paper.

\subsection{$\Pform$ cancellation by block orthogonality}
\label{sec:pchannel:Pform}

$\Pform$ is defined as the form-Hodge contraction of the curvature 2-form.
In components,
\begin{equation}
  \Pform = \frac{1}{4}\sum_{a,b,c,d} R^{ab\,cd}\,(*R_{ab})_{cd},
  \qquad
  (*R_{ab})_{cd} = \frac{1}{2}\epsilon_{cd\,ef}\,R_{ab}{}^{ef},
\end{equation}
with the Hodge dual taken in the Lorentzian metric-aware form-index sense.

\begin{lemma}[Block-orthogonality cancellation]\label{lem:pform}
In the Lorentzian isotropic EC+NY reduced sector studied here, for each
active torsion branch $\AX$, $\VT$, $\MX$ and for each topology $\Sthree$,
$\Tthree$, $\Nilthree$, $\Solthree$, the form-Hodge Pontryagin density
satisfies
\begin{equation}
  \Pform = 0
\end{equation}
identically.
\end{lemma}

\begin{proof}[Proof sketch]
Under the reduced coframe $e^{0}=N\,dt$, $e^{i}=q\,\sigma^{i}$, the
non-zero curvature 2-forms split, for each internal pair, into a
lapse-space block and a purely spatial block.  The form-Hodge dual maps
each block to its complement, but no internal pair admits a contractible
partner, so the pairwise contraction vanishes.  Symbolic verification of all
branches is recorded in Appendix~\ref{app:B}.  This cancellation is a
structural result of the isotropic reduced setting considered here, and
does not hold in general in the full theory or under anisotropic
extensions.  Moreover, $\Pform=0$ is not a claim of dynamical stability.
\end{proof}

\subsection{$\PintMX$ as internal-pair diagnostic}
\label{sec:pchannel:Pint}

On $\AX$ and on $\VT$, the off-shell expression for $\Pint$ vanishes
identically.  On $\MX$ an off-shell non-zero channel appears, but it
returns to zero on the auxiliary shell ($\eta=V=0$ under the real-branch
condition).  This off-shell behaviour does not violate the Hamiltonian
admissibility.

On the $\MX$ branch, the off-shell $\PintMX$ takes the common functional
template across the four topologies,
\begin{equation}
  \PintMX
  = \frac{2V\eta\!\left(-V^{2}q^{2}+9\eta^{2}+\Ctop\right)}{9q^{3}},
\end{equation}
and the topology dependence is localised entirely in the single coefficient
$\Ctop$:

\begin{table}[H]
\centering
\begin{tabular}{lrrr}
\toprule
Topology & $\chi$ & $\Ctop$ & $\Ctop+9\chi$ \\
\midrule
$\Sthree$   & $4$     & $-36$  & $0$ \\
$\Tthree$   & $0$     & $0$    & $0$ \\
$\Nilthree$ & $-1/12$ & $3/4$  & $0$ \\
$\Solthree$ & $-1/3$  & $3$    & $0$ \\
\bottomrule
\end{tabular}
\caption{$\Ctop = -9\chi$ bridge table.}
\label{tab:ctop_bridge}
\end{table}

The coefficient $\Ctop$ is extracted from the raw symbolic expression of
$\PintMX$ by template fitting, without reference to $\chi$.  The scalar
$\chi$ is obtained independently from the $\EH$ / spatial-curvature sector
(see Section~\ref{sec:setup:chi}).  The agreement of these two quantities
is the substance of the $\chi$-bridge result discussed in
Section~\ref{sec:bridge}.

\subsection{$\QNY$ as endpoint channel}
\label{sec:pchannel:QNY}

The Nieh--Yan primitive~\cite{NiehYan1982,ChandiaZanelli1997} is
$\QNY = e^{a}\wedge T_{a}$, and we treat it as an exact-form endpoint
quantity independent of $\Pform$ and $\Pint$.  The branches $\AX$ and
$\MX$ carry an active axial torsion $\eta$, and the reduced pullback reads
\begin{equation}
  \QNY = \frac{6\eta}{q}\,e^{123},
  \qquad
  \int_{\topo} \QNY = 6\Vol_{\topo}\, q^{2}\eta,
\end{equation}
with orientation fixed by $Q(t_{f})-Q(t_{i})$.  Here $\Vol_{\topo}$ is the
unit-coframe volume normalisation factor, which is distinguished from the
torsion variable $V(t)$.  On the $\VT$ branch, the EOM imposes
$\eta=0$, so $\QNY=0$ already off-shell.  On the auxiliary shell, $\AX$ is
on-shell trivial and $\MX$ is boundary-conditional.  Thus $\QNY$ is not a
bulk dynamical mechanism but should be treated as endpoint / boundary
bookkeeping.
